% English abstract -- Mémoire d'Ingénieur d'État.
% One paragraph of context, one of what was done, one of what was found.
% Numbers here must trace to results/*.csv; see the thesis-writing skill.
\thispagestyle{empty}
\vspace*{4cm}
\begin{center}{\Huge\bfseries Abstract}\end{center}
\vskip 1cm
\noindent A drone swarm commanded by voice must act on a stopping command faster than on any command
it has to interpret, and it must do so without depending on infrastructure the flight site may not
provide. This thesis designs, builds and measures a speech-to-swarm runtime whose models all execute
on a Raspberry~Pi~5, and asks whether it meets a latency budget compatible with interactive control,
and whether recognition degrades safely under propeller noise while the controller holds formation.

\vspace{0.3cm}
\noindent The runtime divides one audio stream between two paths. On the reflex path, a two-class
keyword spotter maps \emph{swarm hold} and \emph{swarm abort} directly to a hold or an abort and
preempts language-model inference; a membership rule derived from the cost of a false accept and the
value of the time saved admits exactly these two classes and no movement command. On the parse path,
each utterance is endpointed, transcribed and parsed by a grammar-constrained Qwen2.5-0.5B model. A
semantic validator and a flight state machine check every command, and a rejection never becomes
motion: in flight it becomes a hold. A formation controller with a geometric separation clamp at the
integrator flies five simulated drones.

\vspace{0.3cm}
\noindent The reflex path cancelled 78 of 78 in-flight decodes with no cross-trigger between its
classes, but its p95 latency, 545~ms idle and 547~ms under load, misses the 150~ms budget after the
keyword offset, in the spotter's decision delay, the frame quantisation and the anchor's error,
rather than in computation. The parse path reaches 3{,}122~ms at
p95 against 2{,}500~ms, the shortfall located in speech recognition and in an uncached prompt
prefill. Command recognition is 0.690 on clean audio and 0.590 at 10~dB, below its targets of 0.80
and 0.65, and only 2 of 62 clean-audio failures resolve to a hold or to \texttt{unknown}. Zero
collisions were observed across 150 formation trials, backed by the clamp, which intervened 432
times. An end-to-end run with networking disabled and the live demonstration are left to the
defence.

\vspace{1cm}
\providecommand{\keywords}[1]{\textbf{\textit{Keywords---}} #1}
\keywords{dual-path runtime; keyword spotting; preemption; offline speech recognition;
safety architecture; drone swarm; formation control; Raspberry Pi 5.}
\newpage
